

For \emph{Proof of Reflection} to work, we must alter the \emph{fork rule} (which compares the chain-weight of two candidate best-blocks).
We sum block-weights to get the chain-weight.
PoW chains use \emph{the expected number of hashes to produce a block} as the measure of \emph{work}.\footnote{Alternatively, something that's \emph{linearly convertible} to number-of-hashes (like difficulty) can be used instead.}
So, for mutually reflecting PoW chains, we need a way to \emph{compare and convert} hashes done on L to hashes done on R.

If we are to convert work between PoW chains, then we \emph{must} be able to convert between L-hashes and R-hashes \emph{in a consistent way}.
Consistency, here, means the conversion shouldn't change the outcome of the fork rule; if the fork rule says $L_{i \cdots j,a} > L_{i \cdots j,b}$ in terms of L-hashes, then it should give the same result for \emph{all} cases when the units are changed to R-hashes.

\todoDraftOnly{QuickCheck property tests on the properties described}

% footnote on 'easiest': \footnote{This might be the \emph{only} way to keep things consistent; but for our purposes, it doesn't matter.}

Since both weight and hashes can be \emph{summed}, an easy way for this to work is for us to use a \emph{linear} conversion method (i.e., the units are \emph{proportional}, we have some \emph{constant of conversion} for that specific situation, and all units \emph{share an origin}).
In effect, for any two values, the \emph{ratio} between linearly converted values represented in same units is \emph{constant}, regardless of which units those values are represented in.
This doesn't imply that \emph{mid-states} are linearly convertible, though.
However, if some mid-states \emph{are} linearly convertible (to and from hashes), then \emph{the fork rule should work for the absolute value of those units, too}.

When it comes to comparing work, we're going to normalize some factors against time.
These are the factors which \emph{naturally come as rates}, even though they're often measured in absolutes.
For example: instead of measuring work in hashes, we'll immediately convert to \emph{hashes per second};
instead of measuring \emph{block rewards}, we'll measure \emph{income per second} and \emph{inflation rate}.
Even though the act of block production is \emph{discrete}, blockchains have a \emph{block target time}; a way of keeping things consistent \emph{over time} (at least in the short term).
Since a block represents a discrete amount of work, and has a known, measurable target time, we can unify the two via the idea of \emph{rate of work}, e.g., \emph{hashes per second}.
Making this conversion up-front will simplify the work to come, since block production frequencies don't need to be accounted for.

\aside{
    If we can measure the \emph{rate of work} in some direct way (e.g., number of hashes per block) then we can also \emph{scale block rewards according to both the network difficulty and the demonstrated rate of work}.
    This enables dynamic block production (i.e., blocks with a flexible minimum difficulty).
    This would (probably) be disruptive to a traditional blockchain, due to many stale blocks.
    For UT, it becomes relevant with \autoref{sec:dos-and-dags}, and particularly with \autoref{sec:preventing-dos-attacks}.
}

\textbf{Here's the problem:} \emph{universal and generic} conversion between qualitatively different units is impossible.\footnote{
    WRT ideas around conversion and the role of goals and context, I learned them from Elliot Temple and his \href{https://curi.gumroad.com/l/mhtbA}{Critical Fallibilism Course}.
    For relevant background to those ideas: see \href{https://criticalfallibilism.com/multi-factor-decision-making-math}{Multi-Factor Decision Making Math}.
}
We need a specific \emph{goal} and \emph{context} to successfully convert -- in fact, the \emph{goal} is what \emph{enables} us to judge whether it's a successful method of conversion or not.

The goal of this conversion is the only specific thing we know right now: to convert chain-work so that \emph{PoR works}.
The \emph{conversion} shouldn't introduce any vulnerabilities, and it should work within the constraints that we talked about above.

In order to discuss \emph{how} to convert we need to know the \emph{context} that we operate in.
We know some of the context already: the rest of the world, e.g., we know about some ways to attack blockchains.
That's background context.
What specific contexts can \emph{facilitate} conversion?
Here are two:

% continue on to single-root-token-multi-chain conversion method
